% !TEX root = ../main.tex
\section{Introduction}

After observing data, $x$,
a Bayesian agent updates his uncertainty about
an unknown quantity, $\theta \in \Theta$, from
a prior probability, $\P(\theta)$, to
$\P(\theta|x)$. However,
given human limitations \citep{Tversky1974},
it can be challenging to
interpret and communicate
the full entailments of $\P(\theta|x)$.
These challenges can be addressed by
a \textit{probability summary}, that is,
a simplification of $\P(\theta|x)$ that
categorizes propositions about $\theta$ into
alethic modalities. \red{Better describe what is meant by a summary and how it can be used. Contrast to beliefs.}

Formally, a \probsumm\ is a function,
$\phi(H|x)$, which assigns
a modality to each hypothesis, $H$,
after data $x$ is observed.
Here, the hypothesis, $H$, is 
a subset of $\Theta$ and stands for
the proposition that $\theta \in H$.
When $\phi$ assumes values in $\{0,1\}$,
$\phi(H|x) = 1$ and $\phi(H|x) = 0$ 
are interpreted, respectively, as
$H$ to be taken as true and as false.
For instance, the 
\textit{Lockean thesis} \citep{Foley2009}
states that a proposition $H$ is 
taken as true if and only if 
its probability is higher than 
a threshold, $r$, that is,
$\phi(H|x) = \I(\P(\theta \in H|x) > r)$.
In particular, when $r = 0.5$, a proposition 
taken as true when
it is more likely to be true than false.
Such a \summ \ is useful,
for instance, in most USA civil legal cases.
In such a setting, one satisfies
the burden of proof regarding a claim by
preponderance of evidence, that is, by
showing that it is
more likely than not \citep{Carlson2023}.
When such a standard is met,
US Law is enacted as if the claim were true.

More generally, 
a thernary \probsumm, $\phi$,
assumes values in $\{0, \half, 1\}$.
In this context, 
$\phi(H|x) = 1$ and $\phi(H|x) = 0$ represent,
respectively, taking $H$ as true and as false.
Also, $\phi(H|x) = \half$ represents
an absence of conviction on whether 
to take $H$ as true or false.
More formally:

\begin{definition}
 Let $\sigma(\Theta)$ be
 a $\sigma$-field of 
 measurable subsets of $\Theta$.
 A thernary \probsumm\
 is a function
 $\phi: \sigma(\Theta) 
 \rightarrow \{0, \half, 1\}$.
 When there is no ambiguity,
 $\phi$ is simply said to be
 a \summ.
\end{definition}

\begin{definition}
 A \summ\ $\phi$ is a 
 \stdsumm\ if,
 for every $\phi \in \{0,1\}$.
\end{definition}

Some additional properties are usually 
necessary for $\phi$ to be 
an interpretable \summ.
For instance, let $H$ be the hypothesis that 
it will rain tomorrow.
If both $\phi(H|x) = 1$ and $\phi(H^c|x) = 1$, then
one takes as true both
that it will rain and 
that it will not rain tomorrow.
Such an apparent logical contradiction is
hard to interpret.
One way of avoiding these
contradictions is
to require logical conditions on
the \summ.

In order to obtain such conditions,
one can look at the related literature of
agnostic hypothesis tests.
An agnostic hypothesis test is a function:
$\psi: \sigma(\Theta) 
\rightarrow \left\{0,\half,1\right\}$.
In this context, 
$\psi(H|x) = 1$ and $\psi(H|x) = 0$ represent,
respectively, that
$H$ is rejected and that it is accepted.
Also, $\psi(H|x) = \half$ represents that
$H$ remains undecided.
In this sense, $1-\psi$ is
a \summ, that is,
there exist a one-to-one relationship
between these functions.
By adapting the conditions of logical coherence
for agnostic tests \citep{Esteves2016},
one obtains a definition of 
\textit{logical coherence}
for \summs, as described below:

\begin{definition}[Proper \summ]
 \label{def:proper}
 $\phi$ is a proper \summ \
 if $\phi(\Theta|x) = 1$,
 for every $x \in \sX$.
\end{definition}

\begin{definition}[Logical coherence]
 \label{def:coherence}
 A \summ, $\phi$, is 
 $\kappa$-coherent if, 
 for every $x \in \sX$:
 \begin{enumerate}
  \item $\phi$ is proper,
  \item $\phi(H|x) = 1 - \phi(H^c|x)$,
  for every $H \in \sigma(\Theta)$ (invertibility),
  \item $\phi(H|x) \leq \phi(H'|x)$,
  for every $H, H' \in \sigma(\Theta)$
  such that $H \subset H'$ (monotonicity),
  \item If $\phi(H|x) = 1$ for every $H \in \GG$
  and $|\GG| < \kappa$, then $\phi(\cap\GG|x) = 1$
  ($\kappa$-consonance).
 \end{enumerate}
 Let $|\Theta|^{+}$ be 
 the successor cardinal of $|\Theta|$.
 If $\phi$ is $|\Theta|^{+}$-coherent,
 then it is said to be
 \textit{completely-coherent}.
\end{definition}
\Cref{def:coherence} is also similar to 
property (P1) for beliefs in \citet{Leitgeb2013}.
However, since the latter studies solely
the case in which $\Theta$ is finite,
only $\aleph_0$-consonance is discussed.

This paper studies whether logically coherent are coherent representations of uncertainty. 
\Cref{sec:coh_prob} shows that a \summ \ is
logically coherent if and only if
it admits a probability representation, 
that is, there exists a probability function that
agrees qualitatively with the \summ.
Specifically, \stdsumms \ are 
logically coherent if and only if
they are extreme-valued probabilities.

\red{From another perspective, one might also wish for $\phi$ to be an optimal decision. That is,
after choosing a loss function $L(\phi, \theta)$,
one might wish to choose
$\phi = \arg\min_{\phi^*} \E[L(\phi^*,\theta)]$.
This paper studies the relation
between logical coherence and optimal decisions
in the context of hypothesis tests. This perspective is studied in Section 3.}
